\documentclass{scrartcl}
\usepackage[utf8]{inputenc}
\usepackage{fancyhdr}
\usepackage{graphicx}
\usepackage{dirtytalk}
\usepackage[portrait, margin=1in]{geometry}
\usepackage{amsmath, amssymb}
\setlength{\parskip}{1em}
\setlength{\parindent}{0cm} 
\usepackage[mdthm,simplethm]{test}
\usepackage[utf8]{inputenc}
\usepackage{float}
\usepackage{physics}
\setlength{\parindent}{0pt}

\title{Programs as Data}
\author{albert ye}
\date{May 2022}

\begin{document}
\maketitle

\section{Programming Languages}
Machine language is just 0's and 1's. Assembly has some simpler commands. HL languages have more abstraction
than Assembly and can be interpreted / compiled into a lower level language.

A compiler eventually turns source code into machine code, which is able to be output.
An interpreter turns source code into output. Some implementations, such as CPython, have both a compilation
step (C) and an interpretation step (Python).

\subsection{Phases of an interpreter/compiler}
To comprehend source code, you need lexing, parsing, and an abstract syntax tree (whatever that is).

\vocab{Parsing} involves turning a string representation of an expression into a structured object representing the expression.
For example in Lisp-like languages you take each individual expression (including like \texttt{(} and \texttt{)} and try to see 
if the expression compiles.

\vocab{Lexical analysis} checks each of the tokens and makes sure they aren't malformed and determines their type iteratively.

\vocab{Syntactic analysis} is a recursive process that tries to balance the parentheses and return a tree structure. Symbols will
be the base case, while the combinations of subexpressions will be the recursion.

With this recursive quality, we can make a Scheme compiler in Python (although that's kinda wack)

\end{document}